\documentclass{article}
\usepackage{graphicx}      % Required for inserting images
\usepackage{amsmath}       % Required for math formatting
\usepackage{algorithm}     % Required for the algorithm float environment
\usepackage{algpseudocode} % Required for the algorithmic logic block
\usepackage{booktabs}      % Required for professional table lines
\usepackage{amssymb}

\title{paper2}
\author{hahv}
\date{March 2026}

\begin{document}

	\maketitle

	\section{Introduction}

	To select the optimal hyperparameters for the proposed skip module, we employ a constrained multi-objective optimization procedure on the validation set. Let $R_{\text{base}}$ denote the end-to-end recall and $\mathrm{FAR}_{\text{base}}$ denote the false alarm rate of the baseline pipeline without skipping. For each candidate parameter set $\theta \in \Theta$ obtained by grid search, we evaluate the full pipeline $\text{Read} \rightarrow \text{Skip}(\theta) \rightarrow [\text{DL}]$ and compute three metrics: recall $R(\theta)$, false alarm rate $\mathrm{FAR}(\theta)$, and negative-frame skip ratio $S(\theta)$.

	The negative-frame skip ratio is defined as
	\[
	S(\theta)=
	\frac{\#\ \text{correctly skipped negative frames}}
	{\#\ \text{total negative frames}}.
	\]

	To combine multiple objectives in a single score, the metrics should be expressed on comparable scales. Accordingly, we define the normalized false alarm reduction as
	\[
	\Delta \widetilde{\mathrm{FAR}}(\theta)=
	\max\left(
	0,\,
	\frac{\mathrm{FAR}_{\text{base}}-\mathrm{FAR}(\theta)}
	{\mathrm{FAR}_{\text{base}}}
	\right),
	\]
	which measures the relative reduction in false alarm rate with respect to the baseline system. Under the conservative-gating assumption of the proposed pipeline, $\mathrm{FAR}(\theta)\leq \mathrm{FAR}_{\text{base}}$ should hold theoretically; the $\max(0,\cdot)$ form is retained for robustness and implementation safety.

	To reflect recall preservation in a simpler and more interpretable way, we define the recall-retention term as
	\[
	\widetilde{R}(\theta)=
	1-\frac{R_{\text{base}}-R(\theta)}{\delta_R},
	\]
	where $\delta_R$ is the maximum allowable absolute recall drop. This term equals $1$ when the candidate matches the baseline recall and equals $0$ when the candidate reaches the lowest acceptable recall level $R_{\text{base}}-\delta_R$.

	The optimal parameter set $\theta^*$ is selected as
	\[
	\theta^*=
	\arg\max_{\theta \in \Theta}
	\left[
	w_S S(\theta)
	+
	w_F \Delta \widetilde{\mathrm{FAR}}(\theta)
	+
	w_R \widetilde{R}(\theta)
	\right]
	\quad
	\text{subject to}
	\quad
	R(\theta)\geq R_{\text{base}}-\delta_R,
	\]
	where $\delta_R=0.01$ denotes a 1\% absolute recall-drop tolerance and the nonnegative weights satisfy
	\[
	w_S+w_F+w_R=1.
	\]

	In this work, we set
	\[
	w_S=0.60,\qquad
	w_F=0.20,\qquad
	w_R=0.20,
	\]
	so that skip ratio remains the primary efficiency objective, while false alarm reduction and recall retention are treated as secondary but still explicit preferences. This formulation preserves the safety-first role of recall through the hard constraint, while avoiding the drawback of selecting among feasible candidates using skip ratio alone.

	\textbf{Theoretical justification.}
	The skip module acts as a conservative gate: skipped frames are forced to output ``negative,'' while passed frames are processed by the same downstream detector as in the baseline system. Therefore, the false alarms of the skip-enabled system form a subset of those of the baseline system, implying $\mathrm{FAR}(\theta)\leq \mathrm{FAR}_{\text{base}}$ under the assumed architecture. The main safety risk is thus recall degradation due to false skips, which motivates using recall as both a hard feasibility condition and a soft ranking term.

	\begin{algorithm}[htbp]
		\caption{Skip Module Hyperparameter Selection}
		\begin{algorithmic}[1]
			\Require parameter space $\Theta$, validation set $D_{\text{val}}$, detector $M$, recall tolerance $\delta_R$, weights $w_S,w_F,w_R$
			\Ensure optimal parameters $\theta^*$

			\State baseline $\gets$ evaluate\_pipeline($D_{\text{val}}$, skip=None, $M$)
			\State extract $R_{\text{base}}, \mathrm{FAR}_{\text{base}}$ from baseline
			\State best\_score $\gets -\infty$
			\State best\_$\theta \gets$ None

			\For{$\theta \in \Theta$}
			\State skip $\gets$ SkipModule($\theta$)
			\State results $\gets$ evaluate\_pipeline($D_{\text{val}}$, skip, $M$)
			\State extract $R(\theta)$, $\mathrm{FAR}(\theta)$, $S(\theta)$ from results

			\If{$R(\theta) \geq R_{\text{base}}-\delta_R$}
			\State $\Delta \widetilde{\mathrm{FAR}}(\theta)\gets
			\max\left(0,\frac{\mathrm{FAR}_{\text{base}}-\mathrm{FAR}(\theta)}{\mathrm{FAR}_{\text{base}}}\right)$
			\State $\widetilde{R}(\theta)\gets
			1-\frac{R_{\text{base}}-R(\theta)}{\delta_R}$
			\State score $\gets
			w_S S(\theta)+w_F \Delta \widetilde{\mathrm{FAR}}(\theta)+w_R \widetilde{R}(\theta)$

			\If{score $>$ best\_score}
			\State best\_score $\gets$ score
			\State best\_$\theta \gets \theta$
			\EndIf
			\EndIf
			\EndFor

			\State \Return best\_$\theta$
		\end{algorithmic}
	\end{algorithm}

	\section{Results}

	Table~\ref{tab:skip-selection} shows an example of the validation-time ranking. The selected configuration $\theta_1^*$ satisfies the recall constraint and achieves the highest composite score by jointly balancing skip ratio, false alarm reduction, and recall retention.

	\begin{table}[htbp]
		\centering
		\caption{Example hyperparameter selection results}
		\label{tab:skip-selection}
		\begin{tabular}{lccccc}
			\toprule
			Param Set & Recall Drop & FAR Red. & Skip Ratio & Score & Selected \\
			\midrule
			$\theta_1^*$ & 0.4\% & 42\% & 87.3\% & \textbf{0.728} & $\checkmark$ \\
			$\theta_2$   & 0.8\% & 45\% & 84.1\% & 0.635 & $\times$ \\
			$\theta_3$   & 1.2\% & 51\% & 82.5\% & -- & $\times$ \\
			\bottomrule
		\end{tabular}
	\end{table}

	This formulation is systematic, interpretable, and aligned with the intended role of the skip module in real-time fire/smoke detection: preserve recall first, then prefer candidates that skip more negative frames while still improving operational false alarm behavior.

\end{document}
